\documentclass[12pt,oneside,a4paper,english,brazil]{abntex2}

\usepackage[utf8]{inputenc}
\usepackage[T1]{fontenc}
\usepackage[brazil]{babel}
\usepackage{mathptmx}
\usepackage{amsmath,amssymb}
\usepackage{booktabs}
\usepackage{array}
\usepackage{tabularx}
\usepackage{longtable}
\usepackage{float}
\usepackage{graphicx}
\usepackage{setspace}
\usepackage{xcolor}
\usepackage{colortbl}
\usepackage{enumitem}
\usepackage{listings}
\usepackage{caption}
\usepackage[a4paper,top=3cm,bottom=2cm,left=3cm,right=2cm]{geometry}
\hypersetup{hidelinks}

% Formatação geral conforme ABNT NBR 14724
\setlength{\parindent}{1.25cm}
\setlength{\parskip}{0pt}
\OnehalfSpacing
\renewcommand{\arraystretch}{1.15}
\setlength{\beforechapskip}{-10pt}
\setlength{\afterchapskip}{18pt}

\definecolor{lilasescuro}{RGB}{112,78,150}
\definecolor{lilasclaro}{RGB}{238,230,248}
\definecolor{cinzaclaro}{RGB}{245,245,245}

% Títulos primários e secundários
\renewcommand{\ABNTEXchapterfont}{\normalfont\bfseries\color{lilasescuro}}
\renewcommand{\ABNTEXchapterfontsize}{\fontsize{14}{17}\selectfont}
\renewcommand{\ABNTEXsectionfont}{\normalfont\bfseries}
\renewcommand{\ABNTEXsectionfontsize}{\fontsize{12.5}{15}\selectfont}
\renewcommand{\ABNTEXsubsectionfont}{\normalfont\bfseries}
\renewcommand{\ABNTEXsubsectionfontsize}{\normalsize}

% Paginação no canto superior direito a partir da parte textual
\makepagestyle{abntmain}
\makeoddhead{abntmain}{}{}{\thepage}
\makeevenhead{abntmain}{\thepage}{}{}
\makeoddfoot{abntmain}{}{}{}
\makeevenfoot{abntmain}{}{}{}

\captionsetup{font=small,labelfont=bf,justification=centering,singlelinecheck=false}

\lstset{
  basicstyle=\ttfamily\footnotesize,
  backgroundcolor=\color{cinzaclaro},
  frame=single,
  breaklines=true,
  columns=fullflexible,
  showstringspaces=false,
  keywordstyle=\bfseries\color{lilasescuro},
  commentstyle=\itshape,
  literate={á}{{\'a}}1 {é}{{\'e}}1 {í}{{\'i}}1 {ó}{{\'o}}1 {ú}{{\'u}}1
           {ã}{{\~a}}1 {õ}{{\~o}}1 {ç}{{\c c}}1 {Á}{{\'A}}1 {É}{{\'E}}1
}

\begin{document}
\pretextual
\pagestyle{empty}

\begin{capa}
\begin{center}
{\large\bfseries UNIVERSIDADE FEDERAL DO CEARÁ\par}
\vspace{0.25cm}
{\large TI0111 -- Statistics for Engineers\par}

\vspace{2.8cm}
{\large Allan Martins Gadelha -- 579087\par}
{\large Clarisse Maria Cabral Montenegro -- 580585\par}
{\large Maria Luisa Fernandes de Mendonça -- 581706\par}
{\large Sara de Sousa Magalhães -- 581189\par}

\vfill
{\Large\bfseries RELATÓRIO FINAL -- HOMEWORK 1\par}
\vspace{0.35cm}
{\large\bfseries Análise Estatística de um Sistema de Compartilhamento de Bicicletas\par}

\vfill
{\small\bfseries\color{lilasescuro} Recursos digitais do projeto\par}
\vspace{0.15cm}
{\small GitHub: \href{https://github.com/SaraMag25/Homework-1}{\texttt{github.com/SaraMag25/Homework-1}}\par}
{\small Google Drive: \href{https://drive.google.com/drive/folders/1jq3fmuYgX59cN8P8pv3AWH5V_lezF6Zc?usp=drive_link}{\texttt{Pasta do projeto no Google Drive}}\par}

\vfill
{\large Fortaleza\par}
{\large 2026\par}
\end{center}
\end{capa}

\begin{folhaderosto}
\begin{center}
{\large Allan Martins Gadelha\par}
{\large Clarisse Maria Cabral Montenegro\par}
{\large Maria Luisa Fernandes de Mendonça\par}
{\large Sara de Sousa Magalhães\par}

\vfill
{\Large\bfseries RELATÓRIO FINAL -- HOMEWORK 1\par}
\vspace{0.35cm}
{\large\bfseries Análise Estatística de um Sistema de Compartilhamento de Bicicletas\par}

\vfill
\hspace{0.45\textwidth}
\begin{minipage}{0.50\textwidth}
\SingleSpacing
Relatório final apresentado à disciplina TI0111 -- Statistics for Engineers, da Universidade Federal do Ceará, como parte das atividades avaliativas do Homework 1.
\end{minipage}

\vfill
{\large Fortaleza\par}
{\large 2026\par}
\end{center}
\end{folhaderosto}

\begin{resumo}
Este relatório apresenta a construção, exploração e análise estatística de uma amostra com 300 observações diárias de um sistema de compartilhamento de bicicletas. O trabalho abrange a definição da amostra, validação manual de medidas estatísticas, análise descritiva, identificação de baixa utilização, comparação por estação e condição meteorológica, estudo da relação entre temperatura e demanda, análise temporal e aprofundamentos com regressão e análise de variância. Os procedimentos foram desenvolvidos em R, com apoio de pacotes para manipulação de dados, visualização e leitura e escrita de planilhas, e documentados em LaTeX. Os resultados indicam associação positiva entre temperatura e utilização, forte concentração de dias de baixa utilização no inverno e redução do uso sob condições meteorológicas desfavoráveis.

\vspace{\onelineskip}
\noindent\textbf{Palavras-chave}: estatística descritiva. R. compartilhamento de bicicletas. análise exploratória. série temporal.
\end{resumo}

\pdfbookmark[0]{\contentsname}{toc}
\tableofcontents*
\cleardoublepage

\pdfbookmark[0]{\listfigurename}{lof}
\listoffigures*
\cleardoublepage

\pdfbookmark[0]{\listtablename}{lot}
\listoftables*
\cleardoublepage

\textual
\pagestyle{abntmain}
\aliaspagestyle{chapter}{abntmain}
\chapter{Introdução}

Este relatório consolida o desenvolvimento completo do Homework 1 da disciplina TI0111 -- Statistics for Engineers. O trabalho teve como objetivo aplicar conceitos de estatística descritiva e análise exploratória a um conjunto de dados real referente ao uso diário de um sistema de compartilhamento de bicicletas em uma cidade dos Estados Unidos.

O conjunto original possui 731 observações diárias. A partir da maior matrícula da equipe, $M=581706$, foi calculado o índice inicial da amostra:
\begin{equation}
r = 1 + (M \bmod 100)=7.
\end{equation}
Dessa forma, foram selecionadas 300 observações consecutivas, dos índices 7 a 306, correspondentes ao período de 07/01/2011 a 02/11/2011. Todas as análises subsequentes foram realizadas sobre essa mesma amostra, preservando sua ordem original.

O projeto foi desenvolvido de forma modular, com scripts em R, relatórios em LaTeX, arquivos intermediários em Excel e gráficos gerados em PNG. O versionamento e a colaboração foram organizados por Git e GitHub, enquanto RStudio e VS Code foram utilizados na execução e desenvolvimento dos scripts. A redação final e a organização acadêmica foram conduzidas em LaTeX/Overleaf.

\chapter{Organização dos dados e construção da amostra -- Questão 1}

\section{Estrutura da amostra}

A amostra contém as variáveis \texttt{instant}, \texttt{dteday}, \texttt{season}, \texttt{weathersit}, \texttt{temp}, \texttt{casual} e \texttt{registered}. A variável derivada \texttt{total\_user} foi criada posteriormente pela soma
\begin{equation}
\texttt{total\_user}=\texttt{casual}+\texttt{registered}.
\end{equation}

No R, a base foi estruturada com \texttt{read.csv()}, convertida para os tipos adequados com funções como \texttt{as.Date()}, \texttt{as.integer()} e \texttt{as.numeric()}, e validada com \texttt{stopifnot()}. Também foi utilizado o pacote \texttt{openxlsx} para exportação em formato Excel.

\begin{table}[H]
\centering
\caption{Resumo da amostra específica do grupo.}
\begin{tabular}{lc}
\toprule
\rowcolor{lilasclaro}
\textbf{Informação} & \textbf{Resultado} \\
\midrule
Maior matrícula ($M$) & 581706 \\
Índice inicial ($r$) & 7 \\
Primeiro índice selecionado & 7 \\
Último índice selecionado & 306 \\
Número de observações & 300 \\
Primeira data & 07/01/2011 \\
Última data & 02/11/2011 \\
\bottomrule
\end{tabular}
\end{table}

\section{Verificação manual com as primeiras 10 observações}

Como etapa de validação, as primeiras dez observações foram utilizadas para calcular manualmente média, mediana, moda, variância, desvio padrão e quartis de \texttt{total\_user}. Os principais resultados foram:

\begin{table}[H]
\centering
\caption{Resultados manuais e verificação no R para as dez primeiras observações.}
\begin{tabular}{lcc}
\toprule
\rowcolor{lilasclaro}
\textbf{Medida} & \textbf{Cálculo manual} & \textbf{R} \\
\midrule
Média & 1231,60 & 1231,60 \\
Mediana & 1255,50 & 1255,50 \\
Variância amostral & 44525,60 & 44525,60 \\
Desvio padrão amostral & 211,01 & 211,01 \\
$Q_1$ & 1111,25 & 1172,50 \\
$Q_2$ & 1255,50 & 1255,50 \\
$Q_3$ & 1409,75 & 1384,75 \\
IQR & 298,50 & 212,25 \\
\bottomrule
\end{tabular}
\end{table}

As medidas de tendência central e dispersão coincidiram entre o cálculo manual e o R. A diferença nos quartis decorre do método de interpolação: o cálculo manual utilizou posições baseadas em $(n+1)$, enquanto o R emprega por padrão o método \texttt{type = 7} em \texttt{quantile()}. A partir desse ponto, o método do R foi adotado como padrão para manter consistência ao longo do projeto.

\chapter{Caracterização e análise descritiva -- Questão 2}

\section{Classificação das variáveis e integridade dos dados}

A análise inicial confirmou a ausência de valores faltantes nas 300 observações. As variáveis \texttt{season} e \texttt{weathersit} foram tratadas como categóricas; \texttt{temp} como quantitativa contínua; e \texttt{casual}, \texttt{registered} e \texttt{total\_user} como contagens discretas. A variável \texttt{instant} funciona como identificador sequencial e \texttt{dteday} como referência temporal.

Foram observados 73 dias de inverno, 92 de primavera, 94 de verão e 41 de outono. Quanto às condições meteorológicas, a amostra contém 184 dias de céu limpo, 106 dias nublados e 10 dias de chuva fraca; não houve observações classificadas como chuva forte no recorte utilizado.

\section{Medidas de tendência central}

\begin{table}[H]
\centering
\caption{Medidas de tendência central nas 300 observações.}
\begin{tabular}{lrrp{5.2cm}}
\toprule
\rowcolor{lilasclaro}
\textbf{Variável} & \textbf{Média} & \textbf{Mediana} & \textbf{Moda} \\
\midrule
\texttt{temp} & 21,15 & 22,2 & 26,0 (6 ocorrências) \\
\texttt{casual} & 743,77 & 673 & 639 (3 ocorrências) \\
\texttt{registered} & 2769,45 & 3093 & 16 valores empatados, 2 ocorrências cada \\
\texttt{total\_user} & 3513,21 & 3996 & 9 valores empatados, 2 ocorrências cada \\
\bottomrule
\end{tabular}
\end{table}

Os usuários registrados constituem a maior parcela da utilização do sistema. Para \texttt{total\_user}, a mediana superior à média indica que dias de utilização mais baixa reduzem a média global. A moda possui baixa capacidade descritiva para \texttt{registered} e \texttt{total\_user}, já que diversos valores empatam com frequências pequenas.

\section{Quartis, IQR e valores atípicos}

Para \texttt{total\_user}, foram obtidos:
\begin{equation}
Q_1=2105{,}5,\qquad Q_2=3996,\qquad Q_3=4681,
\end{equation}
\begin{equation}
IQR=2575{,}5.
\end{equation}
Pela regra de $1{,}5\times IQR$, os limites foram
\begin{equation}
L_{inf}=-1757{,}75,\qquad L_{sup}=8544{,}25.
\end{equation}
Como o mínimo da amostra é 431 e o máximo é 6043, nenhuma das 300 observações foi classificada como outlier por esse critério global.

\section{Histograma e boxplot}

\begin{figure}[H]
\centering
\includegraphics[width=0.78\textwidth]{2.4.histogramaq2.png}
\caption{Histograma da distribuição de \texttt{total\_user}.}
\end{figure}

O histograma evidencia concentrações de observações em faixas distintas de utilização. Essa estrutura sugere diferentes regimes de demanda ao longo do período, mas o gráfico isoladamente não permite determinar suas causas.

\begin{figure}[H]
\centering
\includegraphics[width=0.68\textwidth]{2.4.boxplotq2.png}
\caption{Boxplot global de \texttt{total\_user}.}
\end{figure}

O boxplot confirma a posição dos quartis e a ausência de pontos além dos limites definidos pela regra de $1{,}5\,IQR$.

\section{Definição de baixa utilização}

Foi criada a variável binária \texttt{low\_usage} com base em $Q_1$:
\begin{equation}
\texttt{low\_usage}_i=
\begin{cases}
1,&\texttt{total\_user}_i<2105{,}5,\\
0,&\texttt{total\_user}_i\geq2105{,}5.
\end{cases}
\end{equation}
Foram classificados 75 dias como baixa utilização e 225 como demais dias, correspondendo exatamente a 25\% e 75\% da amostra, respectivamente.

\chapter{Associações com estação, clima e temperatura -- Questão 3}

\section{Utilização por estação do ano}

\begin{table}[H]
\centering
\small
\caption{Estatísticas de utilização por estação.}
\begin{tabular}{lrrrrrr}
\toprule
\rowcolor{lilasclaro}
\textbf{Estação} & $n$ & \textbf{Média} & \textbf{Mediana} & \textbf{DP} & \textbf{Dias low} & \textbf{Proporção} \\
\midrule
Inverno & 73 & 1639,82 & 1605,0 & 575,69 & 60 & 82,19\% \\
Primavera & 92 & 3775,17 & 4085,5 & 1138,90 & 11 & 11,96\% \\
Verão & 94 & 4464,36 & 4615,5 & 798,35 & 3 & 3,19\% \\
Outono & 41 & 4080,27 & 4304,0 & 1054,78 & 1 & 2,44\% \\
\bottomrule
\end{tabular}
\end{table}

O inverno apresentou os menores níveis de utilização e concentrou 60 dos 75 dias de baixa utilização. O verão teve a maior média e mediana, enquanto a primavera apresentou o maior desvio padrão absoluto.

\begin{figure}[H]
\centering
\includegraphics[width=0.78\textwidth]{boxplot_estacoes.png}
\caption{Distribuição do total diário de usuários por estação do ano.}
\end{figure}

\section{Condições meteorológicas}

\begin{table}[H]
\centering
\caption{Utilização por condição meteorológica.}
\begin{tabular}{lrrrrr}
\toprule
\rowcolor{lilasclaro}
\textbf{Condição} & \textbf{Média} & \textbf{DP} & \textbf{Dias} & \textbf{Dias low} & \textbf{Proporção} \\
\midrule
Céu limpo & 3851 & 1360 & 184 & 33 & 17,9\% \\
Nublado & 3110 & 1328 & 106 & 35 & 33,0\% \\
Chuva fraca & 1562 & 852 & 10 & 7 & 70,0\% \\
\bottomrule
\end{tabular}
\end{table}

A frequência de baixa utilização cresce com a piora das condições meteorológicas. Em dias de chuva fraca, 70\% das observações ficaram abaixo do primeiro quartil global, enquanto em dias de céu limpo essa proporção foi de 17,9\%.

\begin{figure}[H]
\centering
\includegraphics[width=0.72\textwidth]{boxplot_clima.png}
\caption{Distribuição de usuários por condição meteorológica.}
\end{figure}

\begin{figure}[H]
\centering
\includegraphics[width=0.72\textwidth]{proporcao_low_usage.png}
\caption{Proporção de dias de baixa utilização por condição meteorológica.}
\end{figure}

\section{Temperatura e total de usuários}

O coeficiente de correlação de Pearson entre \texttt{temp} e \texttt{total\_user} foi
\begin{equation}
r\approx0{,}806,
\end{equation}
indicando associação linear positiva forte no conjunto analisado. A interpretação é associativa: o coeficiente não estabelece causalidade e não elimina a influência de estação do ano, clima ou outros fatores.

\begin{figure}[H]
\centering
\includegraphics[width=0.75\textwidth]{3.3.grafico_temp_total_user.png}
\caption{Relação entre temperatura e total diário de usuários.}
\end{figure}

O relatório consolidado da Questão 3 destacou temperatura e estação do ano como as duas características mais associadas à utilização. Na implementação posterior da Questão 4, entretanto, o aprofundamento comparativo utilizou temperatura e condições meteorológicas. Este relatório registra essa evolução do foco analítico sem alterar os resultados originais de cada etapa.

\chapter{Aprofundamento conjunto das análises -- Questão 4}

\section{Série temporal}

A série temporal revela crescimento acentuado da utilização ao longo dos primeiros meses, seguido de patamar elevado e maior variabilidade entre agosto e outubro. O menor valor ocorreu em 27/01/2011, com 431 usuários, e o maior em 04/07/2011, com 6043 usuários.

\begin{figure}[H]
\centering
\includegraphics[width=0.95\textwidth]{4.1.serie_temporal_total_user.png}
\caption{Série temporal do total diário de usuários.}
\end{figure}

Entre janeiro e fevereiro predominam valores menores. Março e abril caracterizam uma fase de crescimento, enquanto maio a julho concentram níveis frequentemente superiores a 4000 usuários. Após julho, a utilização permanece elevada, porém com oscilações mais intensas.

\section{Aprofundamento da associação com temperatura}

A correlação de Pearson foi recalculada e manteve o valor aproximado de $0{,}8058$. Uma regressão linear descritiva foi sobreposta ao gráfico de dispersão para evidenciar a tendência geral.

\begin{figure}[H]
\centering
\includegraphics[width=0.76\textwidth]{4.2.grafico_dispersao_temp.png}
\caption{Dispersão entre temperatura e utilização com tendência linear.}
\end{figure}

\section{Condições meteorológicas e ANOVA}

Para comparar o total de usuários entre as categorias de \texttt{weathersit}, foi aplicado um modelo de análise de variância de uma via:
\begin{equation}
\texttt{total\_user}\sim\texttt{weathersit}.
\end{equation}
O teste apresentou $p\approx2{,}13\times10^{-9}$, indicando diferenças estatisticamente significativas entre as médias dos grupos. Esse resultado confirma diferença entre condições meteorológicas, mas não deve ser interpretado isoladamente como prova de causalidade.

\begin{figure}[H]
\centering
\includegraphics[width=0.74\textwidth]{4.3.boxplot_clima.png}
\caption{Comparação do total de usuários entre condições meteorológicas.}
\end{figure}

\section{Temperatura, baixa utilização e comparação entre grupos}

O script final da Questão 4 incorporou uma rotina robusta para interpretar a escala de \texttt{temp}. Caso os valores estejam entre 0 e 1, a temperatura é convertida por $T_C=41T_{norm}$; se já estiverem em intervalo compatível com graus Celsius, os valores são preservados.

A separação por \texttt{low\_usage} mostrou que os dias de baixa utilização apresentam temperatura média de 11,08~$^\circ$C, contra 24,50~$^\circ$C nos demais dias. A correlação entre temperatura e uso foi positiva nos dois grupos, porém mais fraca entre os dias de baixa utilização:

\begin{table}[H]
\centering
\caption{Temperatura e correlação de Pearson por grupo de utilização.}
\begin{tabular}{lrrrr}
\toprule
\rowcolor{lilasclaro}
\textbf{Grupo} & \textbf{Dias} & \textbf{Temp. média} & \textbf{Intervalo de temp.} & $r(T,U)$ \\
\midrule
Baixa utilização & 75 & 11,08 $^\circ$C & 2,4--27,9 $^\circ$C & 0,277 \\
Demais dias & 225 & 24,50 $^\circ$C & 10,8--34,8 $^\circ$C & 0,568 \\
\bottomrule
\end{tabular}
\end{table}

\begin{figure}[H]
\centering
\includegraphics[width=0.86\textwidth]{4.3.dispersao.png}
\caption{Temperatura e utilização, com diferenciação dos dias de baixa utilização.}
\end{figure}

\begin{figure}[H]
\centering
\includegraphics[width=0.88\textwidth]{4.3.paineis.png}
\caption{Comparação visual entre os grupos de baixa utilização e demais dias.}
\end{figure}

A comparação entre grupos deve ser interpretada com cuidado, pois \texttt{low\_usage} é definida a partir da própria variável \texttt{total\_user}. Assim, a separação restringe diretamente o intervalo da variável resposta e pode alterar correlações e inclinações sem representar efeitos causais distintos.

\chapter{Funções, pacotes e recursos adicionais utilizados no R}

Ao longo do projeto foi necessário ir além das operações estatísticas básicas. A Tabela~\ref{tab:funcoesR} resume os principais pacotes e funções efetivamente empregados nos scripts do repositório.

\small
\begin{longtable}{p{3.2cm}p{4.2cm}p{7.1cm}}
\caption{Principais recursos adicionais de R utilizados no projeto.}\label{tab:funcoesR}\\
\toprule
\rowcolor{lilasclaro}
\textbf{Pacote/área} & \textbf{Função ou recurso} & \textbf{Finalidade no projeto} \\
\midrule
\endfirsthead
\toprule
\rowcolor{lilasclaro}
\textbf{Pacote/área} & \textbf{Função ou recurso} & \textbf{Finalidade no projeto} \\
\midrule
\endhead
Base R & \texttt{read.csv()} & Leitura e reconstrução da base original em formato tabular. \\
Base R & \texttt{as.Date()}, \texttt{as.integer()}, \texttt{as.numeric()}, \texttt{as.factor()} & Conversão explícita de tipos para evitar interpretações incorretas durante os cálculos. \\
Base R & \texttt{factor()} & Conversão dos códigos de estação e clima em categorias com rótulos interpretáveis. \\
Base R & \texttt{order()}, \texttt{sort()} & Ordenação cronológica e ordenação numérica para mediana, quartis e séries temporais. \\
Base R & \texttt{stopifnot()}, \texttt{stop()}, \texttt{file.exists()} & Validação defensiva de quantidade de linhas, índices, existência de arquivos e estrutura esperada. \\
Base R & \texttt{names()}, \texttt{setdiff()}, \texttt{length()} & Verificação de colunas essenciais e identificação de campos ausentes. \\
Base R & \texttt{mean()}, \texttt{median()}, \texttt{sd()}, \texttt{IQR()} & Cálculo das medidas descritivas principais. \\
Base R & \texttt{quantile(type=7)} & Cálculo padronizado dos quartis e definição do limiar de baixa utilização. \\
Base R & \texttt{table()}, \texttt{max()}, \texttt{names()} & Implementação da função personalizada \texttt{calcula\_moda()} para detectar moda ou conjunto amodal. \\
Base R & \texttt{ifelse()} & Criação da variável binária \texttt{low\_usage}. \\
Base R & \texttt{sum()}, \texttt{nrow()}, \texttt{anyNA()} & Contagens, validações e quantificação dos dias de baixa utilização. \\
Base R & \texttt{cor(method="pearson")} & Cálculo da associação linear entre temperatura e utilização. \\
Base R & \texttt{lm()} e \texttt{predict()} & Ajuste e projeção de retas de regressão linear descritivas. \\
Base R & \texttt{aov()} e \texttt{summary()} & ANOVA para comparar médias de utilização entre condições meteorológicas. \\
Base R & \texttt{which.max()}, \texttt{which.min()} & Localização das datas de utilização máxima e mínima na série temporal. \\
Base R & \texttt{png()}, \texttt{dev.off()} & Abertura e fechamento de dispositivos gráficos para salvar imagens diretamente em arquivo. \\
Base R & \texttt{plot()}, \texttt{hist()}, \texttt{boxplot()}, \texttt{points()}, \texttt{text()}, \texttt{abline()} & Construção dos gráficos base utilizados em diferentes etapas. \\
Base R & \texttt{par()}, \texttt{legend()}, \texttt{grid()}, \texttt{lines()}, \texttt{mtext()} & Personalização de margens, painéis, legendas, grades e elementos visuais dos gráficos. \\
Base R & \texttt{commandArgs()} & Permitiu ao script final receber arquivo de entrada e pasta de figuras via linha de comando. \\
Base R & \texttt{dir.create()}, \texttt{file.path()} & Criação e composição segura de caminhos para salvar as figuras geradas. \\
Base R & \texttt{range()}, \texttt{seq()}, \texttt{sprintf()} & Definição de limites gráficos, grades numéricas e rótulos formatados. \\
\texttt{readxl} & \texttt{read\_excel()} & Leitura dos arquivos intermediários \texttt{.xlsx} produzidos ao longo das questões. \\
\texttt{openxlsx} & \texttt{write.xlsx()} & Exportação das bases processadas para planilhas reutilizáveis. \\
\texttt{dplyr} & \texttt{mutate()}, \texttt{group\_by()}, \texttt{summarise()}, \texttt{n()} & Criação de novas variáveis e resumos por estação e condição meteorológica. \\
\texttt{dplyr} & \texttt{|>} e \texttt{\%>\%} & Encadeamento das transformações de dados em etapas legíveis. \\
\texttt{ggplot2} & \texttt{ggplot()}, \texttt{aes()}, \texttt{geom\_boxplot()}, \texttt{geom\_col()} & Construção de visualizações por grupo e gráficos de proporções. \\
\texttt{ggplot2} & \texttt{geom\_hline()}, \texttt{labs()}, \texttt{theme\_minimal()}, \texttt{guides()} & Inclusão de limite de $Q_1$, rótulos, tema visual e controle de legendas. \\
\texttt{ggplot2} & \texttt{ggsave()} & Salvamento de gráficos diretamente a partir de scripts executados no terminal/VS Code. \\
\texttt{scales} & \texttt{percent\_format()} & Formatação do eixo de proporções em percentual. \\
\bottomrule
\end{longtable}
\normalsize

\section{Função personalizada para a moda}

A necessidade de lidar com distribuições amodais ou multimodais levou à criação de uma função própria. Ela conta as frequências com \texttt{table()}, encontra a maior frequência e devolve todos os valores empatados ou a indicação de conjunto amodal:

\begin{lstlisting}[language=R]
calcula_moda <- function(v) {
  frequencias <- table(v)
  maior_frequencia <- max(frequencias)

  if (maior_frequencia == 1) {
    return("Amodal")
  } else {
    return(as.numeric(names(frequencias)[frequencias == maior_frequencia]))
  }
}
\end{lstlisting}

\section{Boas práticas de robustez incorporadas}

Além dos cálculos, os scripts incluíram verificações de integridade e rotinas de execução mais robustas. Entre elas, destacam-se a validação da existência de arquivos antes da leitura, a verificação de colunas obrigatórias, o controle do número de observações com \texttt{stopifnot()}, o uso de \texttt{anyNA()} para detectar ausências, a criação automática de diretórios de saída e a detecção da escala de temperatura antes de convertê-la para graus Celsius.

Esses recursos não alteram a teoria estatística utilizada, mas tornam o código mais seguro, reproduzível e reutilizável.

\chapter{Fluxo de trabalho e ferramentas de desenvolvimento}

O projeto foi organizado de forma colaborativa e modular. As principais ferramentas foram:

\begin{itemize}[leftmargin=1.2cm]
\item \textbf{R e RStudio}: execução dos scripts, verificação de objetos, inspeção das bases e produção de gráficos;
\item \textbf{Visual Studio Code}: edição de scripts, execução via terminal e organização local dos arquivos;
\item \textbf{LaTeX e Overleaf}: construção dos relatórios parciais e consolidação do relatório acadêmico;
\item \textbf{Git e GitHub}: controle de versões, histórico de desenvolvimento e integração do trabalho produzido pelos integrantes;
\item \textbf{Trello}: organização das tarefas e acompanhamento do progresso;
\item \textbf{Excel/XLSX}: formato intermediário utilizado para transportar os conjuntos processados entre etapas.
\end{itemize}

A organização por questão permitiu manter scripts, documentos e imagens próximos de suas respectivas análises, reduzindo ambiguidade sobre a origem de cada resultado.


\chapter{Contribuições da equipe}

O trabalho foi desenvolvido de forma colaborativa, com divisão das etapas entre os quatro integrantes e integração contínua pelo GitHub. A síntese abaixo foi organizada a partir das atividades registradas no histórico do repositório e dos arquivos finais produzidos em cada questão. Como houve revisões, reorganizações e integração entre diferentes etapas, as contribuições são apresentadas como principais responsabilidades registradas, sem pressupor exclusividade sobre cada arquivo.

\section{Allan Martins Gadelha -- 579087}

As principais contribuições de Allan concentraram-se no processamento da base e em análises estatísticas e sazonais. Entre as atividades registradas, destacam-se:

\begin{itemize}[leftmargin=1.2cm]
\item desenvolvimento de script em R para processar a base em Excel e incorporar a variável \texttt{total\_user};
\item apoio à preparação do conjunto de dados utilizado nas etapas posteriores;
\item desenvolvimento da análise de utilização por estação do ano, incluindo medidas descritivas e visualizações;
\item produção e inclusão de documentação em LaTeX para partes das análises estatísticas;
\item participação na consolidação da Questão 2 e em análises adicionais relacionadas à temperatura e ao comportamento do sistema.
\end{itemize}

\section{Clarisse Maria Cabral Montenegro -- 580585}

Clarisse atuou principalmente na construção e validação inicial da amostra, em análises descritivas e na série temporal. Suas principais contribuições registradas foram:

\begin{itemize}[leftmargin=1.2cm]
\item construção da tabela da amostra específica do grupo e organização da Questão 1;
\item elaboração e revisão do relatório consolidado da Questão 1;
\item implementação dos cálculos de medidas de tendência central;
\item criação da variável \texttt{low\_usage} a partir do primeiro quartil;
\item desenvolvimento da análise da relação entre temperatura e \texttt{total\_user}, com organização dos arquivos correspondentes;
\item construção da série temporal da Questão 4, incluindo identificação dos valores máximo e mínimo;
\item revisão da estrutura e organização final do repositório.
\end{itemize}

\section{Maria Luisa Fernandes de Mendonça -- 581706}

Maria Luisa contribuiu especialmente para a produção de visualizações, organização dos materiais e consolidação dos relatórios das etapas finais. Entre as principais atividades registradas estão:

\begin{itemize}[leftmargin=1.2cm]
\item desenvolvimento do histograma e do boxplot de \texttt{total\_user} em R e redação da respectiva análise em LaTeX;
\item organização de pastas, imagens e arquivos intermediários das Questões 2 e 3;
\item inclusão e organização das imagens utilizadas na análise da relação entre temperatura e utilização;
\item consolidação do relatório da Questão 3 em LaTeX;
\item organização e integração de partes da Questão 4;
\item elaboração da conclusão da Questão 4 e participação na consolidação do relatório correspondente;
\item preparação e organização do conjunto de imagens finais do projeto.
\end{itemize}

\section{Sara de Sousa Magalhães -- 581189}

Sara teve participação importante na estruturação do repositório, nos cálculos de quartis e nas análises de clima e associação estatística. Suas principais contribuições registradas incluem:

\begin{itemize}[leftmargin=1.2cm]
\item criação e manutenção inicial do repositório, documentação do projeto e atualização do \texttt{README};
\item complementação de partes da Questão 1 e integração dos arquivos iniciais;
\item desenvolvimento do script da Questão 2 para cálculo de quartis, intervalo interquartil e limites de detecção de outliers;
\item redação da documentação em LaTeX correspondente aos quartis;
\item desenvolvimento da análise de condições meteorológicas, incluindo scripts, gráficos e interpretação;
\item desenvolvimento das análises da Questão 4 envolvendo associação entre temperatura, clima e utilização, com scripts R, gráficos e textos em LaTeX;
\item realização de integrações e merges entre as contribuições dos diferentes integrantes e manutenção da organização do repositório.
\end{itemize}

\section{Trabalho colaborativo}

Além das tarefas individuais, o desenvolvimento exigiu integração frequente entre as partes. Os integrantes compartilharam bases intermediárias, revisaram resultados produzidos por colegas, ajustaram caminhos de arquivos, reorganizaram diretórios e consolidaram documentos em LaTeX. O uso do Git e do GitHub permitiu registrar essas etapas e manter um histórico das alterações realizadas ao longo do projeto.

\chapter{Síntese dos principais resultados}

\begin{table}[H]
\centering
\caption{Síntese dos resultados centrais do projeto.}
\begin{tabularx}{\textwidth}{p{4.2cm}X}
\toprule
\rowcolor{lilasclaro}
\textbf{Resultado} & \textbf{Síntese} \\
\midrule
Amostra & 300 dias consecutivos, 07/01/2011 a 02/11/2011. \\
Demanda média & 3513,21 usuários/dia; mediana 3996. \\
Quartis & $Q_1=2105{,}5$, $Q_2=3996$, $Q_3=4681$. \\
Baixa utilização & 75 dias, equivalentes a 25\% da amostra. \\
Outliers globais & Nenhum pela regra de $1{,}5\times IQR$. \\
Estação mais associada à baixa utilização & Inverno: 60 de 73 dias em low usage (82,19\%). \\
Estação com maior média & Verão: 4464,36 usuários/dia. \\
Clima com maior média & Céu limpo: 3851 usuários/dia. \\
Chuva fraca & Média de 1562 usuários/dia e 70\% de dias em low usage. \\
Temperatura x utilização & Correlação de Pearson $r\approx0{,}806$. \\
Série temporal & Mínimo de 431 em 27/01/2011; máximo de 6043 em 04/07/2011. \\
ANOVA por clima & $p\approx2{,}13\times10^{-9}$, indicando diferenças entre médias dos grupos meteorológicos. \\
Correlação por grupo & $r\approx0{,}277$ em low usage e $r\approx0{,}568$ nos demais dias. \\
\bottomrule
\end{tabularx}
\end{table}

\chapter{Conclusão final do projeto}

O desenvolvimento do Homework 1 permitiu percorrer um fluxo completo de análise estatística: seleção e validação da amostra, verificação manual dos cálculos, caracterização das variáveis, construção de medidas descritivas, identificação de quartis, criação de uma variável derivada de baixa utilização, produção de visualizações e investigação de associações entre fatores ambientais e a demanda do sistema.

Os resultados mostram que a utilização não se mantém constante ao longo do período. A demanda é significativamente menor no inverno e em condições meteorológicas desfavoráveis, enquanto os maiores níveis de uso se concentram em períodos mais quentes. A temperatura apresentou forte associação linear positiva com o total de usuários ($r\approx0{,}806$), mas a análise também demonstrou que uma única variável não explica toda a variabilidade observada. A existência de dias quentes com baixa procura e a forte concentração de \texttt{low\_usage} no inverno evidenciam a importância de interpretar conjuntamente temperatura, estação, clima e evolução temporal.

Do ponto de vista computacional, o trabalho também consolidou habilidades de manipulação de dados e programação em R. Além de funções estatísticas clássicas, foram utilizadas rotinas de validação, manipulação de arquivos, programação de funções próprias, agrupamentos com \texttt{dplyr}, gráficos com \texttt{ggplot2} e gráficos base, exportação de planilhas, dispositivos gráficos e execução por linha de comando. Esse conjunto de recursos tornou a análise mais organizada, reprodutível e adequada a um fluxo colaborativo versionado em GitHub.

Por fim, o projeto reforçou a necessidade de distinguir associação de causalidade. Correlações elevadas e diferenças estatisticamente significativas entre grupos fornecem evidências relevantes sobre o comportamento da demanda, mas devem ser interpretadas no contexto do conjunto de dados e das limitações do desenho observacional.


\postextual
\end{document}